\begin{resumo}
O uso de aplicativos móveis tem se tornado cada vez mais indispensável no cotidiano da vidas das pessoas. A existência de um aplicativo para tudo o que fazemos nos torna dependentes desse pequeno aparelho. A área da saúde por ser uma das mais delicadas, demanda a criação de mais aplicativos com o enfoque em cuidados pessoais e gerenciamento de dados. Em vista disso, o objetivo deste trabalho é o desenvolvimento de um aplicativo híbrido para a área da saúde. Para alcançar esse objetivo, foi realizada uma revisão de literatura relacionada visando levantar os principais trabalhos referentes ao gerenciamento de informações pessoais na área da saúde. Adicionalmente, foi elaborado um questionário, para a coleta e análise de dados, assim para entender melhor a necessidade do público alvo. A partir da análise das respostas do questionário foi criado um documento de requisitos, contendo requisitos funcionais e não funcionais, seguido de um diagrama de casos de uso. 


 \vspace{\onelineskip}
    
 \noindent
 \textbf{Palavras-chave}: Aprendizado de Máquina. Linguística Computacional. Sabotagens.
\end{resumo}
